\section{Introduction}
\label{sec:intro}

Large language model agents increasingly alternate between reasoning,
retrieval, and answer generation~\cite{yao2023react,nakano2021webgpt,
trivedi2023ircot,jin2025searchr1}. Each retrieval changes both the cost of an
episode and the evidence available to the answer generator. It may supply a
missing fact, repeat information after the answer has stabilized, or introduce
distracting evidence that changes a correct answer~\cite{chen2024overthinking,
cuadron2025danger,greshake2023notwhat}. Stopping is therefore not merely an
efficiency control. It is part of the answer policy because it sets when
the system commits to its collected evidence.

\begin{figure*}[t]
  \centering
  % Figure 1: overview. Same layout as the accepted figure (stopping loop on top,
% training pipeline below), redrawn in TikZ so all text uses the paper's font.
% The example panel shows HotpotQA id 5a7df1095542990b8f503afe in
% results/measured/run_corpus_guarded_qwen3b.json (prefix EM 0,0,1,0,0;
% learned guard stops after search 3).
% Icons: Microsoft Fluent Emoji (3D), MIT License; see figures/icons/SOURCE.txt.
\definecolor{fInk}{HTML}{2B3441}
\definecolor{fMute}{HTML}{6B7280}
\definecolor{fNavy}{HTML}{2F4B66}
\definecolor{fTeal}{HTML}{2F7470}
\definecolor{fRust}{HTML}{B4583C}
\definecolor{fOchre}{HTML}{A8781F}
\definecolor{fPlum}{HTML}{6A4C6B}
\definecolor{fPaper}{HTML}{F7EFE6}
\definecolor{fSky}{HTML}{EAF0F3}
\definecolor{fCard}{HTML}{FFFDF9}
\providecommand{\ficon}[2][0.42cm]{\raisebox{-0.28\height}{\includegraphics[height=#1]{figures/icons/#2}}}
\begin{tikzpicture}[
  font=\footnotesize, line join=round, x=1cm, y=1cm,
  every node/.append style={execute at begin node={\hyphenpenalty=10000\exhyphenpenalty=10000}},
  band/.style={rounded corners=7pt, line width=0.5pt},
  box/.style={rounded corners=5pt, fill=fCard, draw=#1!70, line width=0.8pt,
    minimum height=0.80cm, inner sep=0pt},
  flow/.style={-{Stealth[length=4.5pt,width=5pt,round]}, line width=0.9pt, draw=fInk!65,
    rounded corners=6pt, line cap=round},
  train/.style={-{Stealth[length=4.5pt,width=5pt,round]}, line width=0.9pt, draw=fRust!85,
    dash pattern=on 3.2pt off 2.2pt, rounded corners=6pt, line cap=round},
  lbl/.style={font=\scriptsize\itshape, text=fMute, inner sep=1.5pt},
  sidelabel/.style={rotate=90, font=\scriptsize\scshape},
  chip/.style={rounded corners=4pt, fill=#1!10, draw=#1!60, line width=0.5pt,
    inner sep=1.5pt, inner xsep=2.4pt, text=fInk, anchor=west, font=\scriptsize},
  tag/.style={rounded corners=3pt, fill=#1, text=white, inner sep=1.6pt,
    inner xsep=2.6pt, anchor=west, font=\scriptsize},
  num/.style={circle, fill=fInk!70, text=white, inner sep=0pt, minimum size=0.30cm,
    font=\tiny\bfseries, anchor=west},
]
% card: one big icon and a short bold title. All cards are 2.60cm wide on a
% 3.00cm pitch, so every gap between neighbouring cards is 0.40cm.
% \fbx{name}{x}{y}{width}{color}{icon}{title}
\newcommand{\fbx}[7]{%
  \node[box=#5, minimum width=#4cm] (#1) at (#2,#3) {};
  \node[anchor=west, inner sep=0pt] (#1i) at ([xshift=0.11cm]#1.west) {\ficon[0.60cm]{#6}};
  \node[anchor=west, inner sep=0pt, align=left, text=#5, font=\footnotesize\bfseries]
    at ([xshift=0.11cm]#1i.east) {#7};
}

% ------------------------------------------------------------- bands
\filldraw[band, fill=fSky, draw=fNavy!25] (0,1.12) rectangle (12.45,3.42);
\node[sidelabel, text=fNavy!80] at (0.25,2.27) {at test time};
\filldraw[band, fill=fPaper, draw=fRust!25] (0,0) rectangle (12.45,0.98);
\node[sidelabel, text=fRust!90] at (0.25,0.49) {training};

% ------------------------------------------------------------- top: the loop
\fbx{st}{1.90}{2.92}{2.60}{fNavy}{brain_3d}{What the\\agent knows}
\fbx{hd}{4.90}{2.92}{2.60}{fPlum}{robot_3d}{Learned\\stop score}
\draw[fPlum, line width=1.6pt, rounded corners=5pt] (hd.south west) rectangle (hd.north east);
\node[diamond, aspect=2.0, fill=fOchre!12, draw=fOchre!80, line width=0.8pt,
      inner sep=0.2pt, align=center, text=fInk] (dc) at (7.90,2.92)
  {\ficon[0.44cm]{vertical_traffic_light_3d}\,\textbf{Stop?}};
\fbx{an}{10.90}{2.92}{2.60}{fTeal}{stop_sign_3d}{Stop and\\answer}

\fbx{env}{7.90}{1.62}{3.05}{fMute}{snowflake_3d}{Frozen LLM\\searches again}
\node[anchor=east, inner sep=0pt] at ([xshift=-0.09cm]env.east)
  {\ficon[0.48cm]{magnifying_glass_tilted_left_3d}};

\draw[flow] (st.east) -- (hd.west);
\draw[flow] (hd.east) -- (dc.west);
\draw[flow] (dc.east) -- node[lbl, above] {yes} (an.west);
\draw[flow] (dc.south) -- node[lbl, right] {no} (env.north);
\draw[flow] (env.west) -- (st.south |- env.west)
  node[lbl, above, pos=0.72] {\ficon[0.32cm]{books_3d}\,new evidence} -- (st.south);

% ------------------------------------------------------------- bottom: training
\fbx{lg}{1.90}{0.49}{2.60}{fRust}{film_frames_3d}{Logged runs}
\fbx{pf}{4.90}{0.49}{2.60}{fRust}{scissors_3d}{Answer after\\every search}
\fbx{tg}{7.90}{0.49}{2.60}{fRust}{bullseye_3d}{Three training\\targets}
\fbx{lk}{10.90}{0.49}{2.60}{fRust}{locked_3d}{Tune, then\\freeze}

\draw[train] (lg.east) -- (pf.west);
\draw[train] (pf.east) -- (tg.west);
\draw[train] (tg.east) -- (lk.west);
\draw[train] (pf.north) -- node[lbl, text=fRust, pos=0.22, right] {train} (hd.south);
\draw[train] (lk.east) -- (12.72,0.49) -- (12.72,2.92) -- (an.east);
\node[fill=white, circle, inner sep=0.6pt] at (12.72,1.70) {\ficon[0.38cm]{snowflake_3d}};

% ------------------------------------------------------------- example panel
\node[rounded corners=6pt, fill=white, draw=fNavy!45, line width=0.6pt,
      dash pattern=on 1.5pt off 1pt, minimum width=3.95cm, minimum height=3.42cm,
      anchor=south west] (ex) at (13.00,0) {};
\node[anchor=north west, inner sep=0pt] at ([shift={(0.10,-0.12)}]ex.north west)
  {\ficon[0.56cm]{thinking_face_3d}};
\node[anchor=north west, text width=3.05cm, align=left, inner sep=0pt, text=fInk,
      font=\scriptsize] at ([shift={(0.76,-0.12)}]ex.north west)
  {\textbf{Example.} \textit{Who was born first, Thomas Carr or Joyce Wieland?}};
\node[anchor=north west, inner sep=0pt, text=fMute, font=\scriptsize\itshape]
  at ([shift={(0.12,-0.84)}]ex.north west) {answer after each search};
\foreach \k/\ans/\mark/\col in {1/Joyce Wieland/cross_mark_3d/fRust,
                               2/Thomas Carr/check_mark_button_3d/fTeal,
                               3/Joyce Wieland/cross_mark_3d/fRust}{
  \node[num] (n\k) at ([shift={(0.12,-0.96-0.40*\k)}]ex.north west) {\k};
  \node[chip=\col] (a\k) at ([xshift=0.07cm]n\k.east)
    {\makebox[1.55cm][l]{\ans}\ficon[0.28cm]{\mark}};
}
\node[anchor=west, inner sep=0pt, font=\scriptsize\bfseries, text=fTeal]
  (r2) at ([xshift=0.10cm]a2.east) {$\leftarrow$\,\ficon[0.28cm]{keycap_2_3d}};
\node[anchor=west, inner sep=0pt, font=\scriptsize\bfseries, text=fRust]
  (r3) at ([xshift=0.10cm]a3.east) {$\leftarrow$\,\ficon[0.28cm]{robot_3d}};
\node[tag=fTeal] at ([shift={(0.12,-2.68)}]ex.north west)
  {\ficon[0.27cm]{keycap_2_3d}\,two fixed searches: right};
\node[tag=fRust] at ([shift={(0.12,-3.08)}]ex.north west)
  {\ficon[0.27cm]{robot_3d}\,learned rule stops: wrong};
\end{tikzpicture}

  \caption{Top: a logged run in which the learned rule stops too late, while a
  fixed two-search budget stops on the right answer. Bottom: the four-step audit.}
  \label{fig:overview}
  \Description{Two-row diagram. The top row follows one logged HotpotQA run:
  after searches one to four the agent answers Joyce Wieland, Thomas Carr
  (correct), Joyce Wieland, and Joyce Wieland. A fixed two-search rule stops at
  the correct answer, while the learned rule stops after the third search with a
  wrong answer. The bottom row shows the audit: record answers and costs after
  every search, learn when to stop, choose and freeze rules on development
  questions, then compare learned and simple rules at equal cost on fresh
  questions.}
\end{figure*}

\paragraph{Motivating example.}
For ``Who was born first, Thomas Carr or Joyce Wieland?'', one logged HotpotQA
trajectory answers Joyce Wieland after one search, Thomas Carr after two searches,
and Joyce Wieland after three. The learned guard stops after that third search
and is wrong. A fixed two-search budget is cheaper and correct. This trace
illustrates why full search alone is an insufficient comparator. We compare
development-selected policies and randomized mixtures on their hull at matched
cost.

The relevant frontier depends on the execution architecture. Our measured
loop generates an answer at every prefix and therefore pays $1+2T$ LLM calls
after $T$ searches. This schedule makes answer repetition a fair low-overhead
baseline because the answers already exist. A fixed budget can instead defer
answer generation until the end and pay $1+T$ calls. We report both the
executed ledger and this answer-once counterfactual. Under the latter, guarded
FBG is dominated in 17 of 18 matched-cost cells and repetition is dominated in
all 18. Cost accounting is therefore part of the controller definition, not a
presentation choice.

Our audit compares four logit-free signals: future-best gain (FBG), binary
future improvement, current-answer error, and gold-support evidence coverage.
FBG labels the largest later token-F1 improvement on one realized continuation,
not a Bellman value. Five-fold chain-level cross-fitting and frozen
development selection precede untouched-question evaluation
(Figure~\ref{fig:overview}). We distinguish point, cellwise, and family-wise
verdicts, decoding variation, and independent question sampling.

Our evaluation spans 26 core cells across four question-answering populations
and nine frozen backbones. Fourteen core cells retain complete prefix-level
cost instrumentation. We combine these cells with four crossed-backbone cells
only for the principal 18-cell matched-cost family. We also freeze a separate
MultiHopRAG news-corpus cell before fitting and test evaluation. Keeping
these populations distinct preserves the provenance, cost records, and
eligible policy families behind each conclusion.

Our findings distinguish stopping opportunity from learned policy value. We
find that FBG-only removes 27 to 57\% of searches on the nine-backbone
HotpotQA panel. Across the 26-cell core, however, a simple policy
Pareto-dominates guarded FBG in 15
cells, while the guard creates a point-frontier improvement over the selected
fixed budget in five. Our hindsight oracle then establishes that the logged
trajectories contain substantial headroom. It improves full search in every
fully instrumented core cell by 4.84 EM points on average. The learned-policy
gap therefore lies in identifying useful exits, not in whether they exist.

The cross-signal analysis makes that diagnosis more precise. FBG orders
stopped-answer risk better than random in 7 of 18 cells, compared with all 18
for our current-error head. Binary future-improvement and error-head policies
produce positive cellwise estimates, yet no policy with an available
simultaneous analysis retains a positive max-t interval. Same-question decoding
reruns preserve the primary non-inferiority decision, whereas independent
question draws produce both passing and failing decisions. Our external news
cell meets its frozen endpoint, while its matched-cost FBG advantage remains
unresolved. These outcomes expose four separate requirements in the deployment
argument: prediction quality, threshold selection, population stability, and
policy utility.

We contribute a selection-aware frontier protocol, a cross-signal empirical
diagnosis, and an audit package with trajectories, prefix costs, selections,
policy cards, and analysis code. These contributions concern shallow per-cell
controllers in four-search replay. Larger training regimes, richer internal
states, and joint query-and-stop policies require separate evaluation.
